\documentclass[12pt,letterpaper]{article}
\usepackage[T1]{fontenc}
\usepackage[utf8]{inputenc}
\usepackage{amsmath,amsthm,mathtools,mathrsfs}
\usepackage{newtxtext,newtxmath}
\usepackage{microtype}
\usepackage[letterpaper,margin=1in]{geometry}
\usepackage{enumitem,booktabs,array,longtable}
\usepackage{xcolor}
\usepackage[hidelinks]{hyperref}
\usepackage{cleveref}

\newtheorem{theorem}{Theorem}[section]
\newtheorem{proposition}[theorem]{Proposition}
\newtheorem{lemma}[theorem]{Lemma}
\newtheorem{conjecture}[theorem]{Conjecture}
\newtheorem{corollary}[theorem]{Corollary}
\theoremstyle{definition}
\newtheorem{definition}[theorem]{Definition}
\newtheorem{example}[theorem]{Example}
\theoremstyle{remark}
\newtheorem{remark}[theorem]{Remark}

\newcommand{\F}{\mathbf{F}}
\newcommand{\C}{\mathbb C}
\newcommand{\R}{\mathbb R}
\newcommand{\Z}{\mathbb Z}
\newcommand{\aaa}{\mathbb{A}}
\newcommand{\bG}{\overline{G}}
\newcommand{\bA}{\overline{A}}
\newcommand{\bC}{\overline{C}}
\newcommand{\bw}{\overline{w}}
\newcommand{\pr}{\mathrm{pr}}
\newcommand{\Pit}{\Pi_t(\bG)}
\newcommand{\tr}{\operatorname{tr}}
\newcommand{\Cent}{\operatorname{Cent}}
\newcommand{\Ind}{\operatorname{Ind}}
\newcommand{\Span}{\operatorname{span}}
\newcommand{\cD}{\mathcal D}
\newcommand{\cS}{\mathcal S}
\newcommand{\calC}{\mathcal C}
\newcommand{\sF}{\mathscr F}
\newcommand{\sJ}{\mathscr J}
\newcommand{\vphi}{\varphi}
\newcommand{\verify}[1]{\textit{[verification pending: #1]}}

\begin{document}

\title{A hitch-hiker's guide to genuine invariant distributions\\
 on metaplectic covers of $SL(2)$ over a $p$-adic field\thanks{Corrected and expanded synthesis of the author's unpublished manuscript dated
 January 4, 1994, and of the $SL(2)$ portions of
 ``Invariant distributions on the $n$-fold metaplectic covers of $GL(r,F)$, $F$ $p$-adic,''
 \emph{Journal of Fourier Analysis and Applications} \textbf{7} (2001), 343--358.
 This draft has not been refereed independently and does not replace the published article.}}
\author{David Joyner\\\texttt{wdjoyner@gmail.com}}
\date{Draft 2: November 21, 2026\\Original notes: January 4, 1994}
\maketitle

\begin{abstract}
This corrected and expanded version of the author's 1994 manuscript was prepared
with the help of AI, specifically ChatGPT (OpenAI) and Claude (Anthropic); see
the Acknowledgments for details.  The paper revisits the original notes on
genuine invariant distributions for
finite central covers of $SL(2,F)$, and incorporates the relevant material from a
subsequent article published in 2001.  We work over a finite extension $F/\mathbb Q_p$
containing $\mu_n$, with $p\nmid 2n$.  After describing the cover, its genuine
representations, and principal-series spectral parameters, we explain the
factorization of invariant distributions through tempered character traces.
The principal algebraic factorization follows from a tempered-character density
statement for covering groups, with the hypotheses and status of the cited
literature made explicit.  A further representation of invariant tempered
distributions as continuous functionals on a spectral Paley--Wiener space is
stated under the precise topological Paley--Wiener input.  Unverified
principal-series image formulas and classification details are identified rather
than treated as independently established theorems.

\end{abstract}


\section{Introduction}\label{sec:intro}
This paper concerns Fourier analysis on a finite central extension of
$SL(2,F)$, where $F$ is a non-Archimedean local field.  On a commutative
group the Fourier transform is parametrized by characters.  On a noncommutative
cover one instead obtains operators $\pi(f)$ and their traces
$\Theta_\pi(f)$, with $\pi$ running through genuine representations.
A \emph{genuine} representation is one on which the central kernel $\mu_n$
acts through its tautological character.  An \emph{invariant distribution}
is a continuous linear functional on compactly supported test functions
that is unchanged by conjugation.  The principal question is how much of
such a distribution is determined by the traces of irreducible tempered
representations.

The starting point is the author's unpublished 1994 manuscript
\cite{Jo94}.  That manuscript described the local cocycle, representations
of the covering torus, principal and complementary series, and proposed
Paley--Wiener descriptions.  It explicitly stated that portions of the
classification and Fourier-image problem were conjectural.  The author's
2001 article \cite{Jo01} subsequently developed an invariant-distribution
argument for $SL(2)$ as well as for certain covers of $GL(r)$.  Here we
consider only the former group.  Removing the general-linear component
allows the local rank-one structure and its limitations to be made explicit.

There are two logically distinct assertions.  First, if a compactly
supported genuine test function has zero trace against every genuine
irreducible tempered representation, a tempered-character density theorem
identifies it with a sum of conjugation differences.  Consequently every
invariant distribution factors algebraically, and uniquely, through the
\emph{image} of the tempered trace map.  Second, describing a \emph{tempered}
invariant distribution as a continuous spectral functional requires a
specified topology on a Paley--Wiener space and a continuous Fourier
surjection with the appropriate kernel.  The second statement does not
follow formally from the first.  See Theorems~\ref{thm:trace}
and \ref{thm:main}.

\begin{example}[Why a spectral functional need not have a density]\label{ex:derivative}
On a one-parameter circle of unitary characters, the rule
$T(h)=h'(0)$ is a distribution, but it is not integration against an
ordinary locally integrable density.  Thus an expression
$\int\Theta_\pi(f)m(\pi)\,d\mu(\pi)$ for a general invariant
distribution must be interpreted as the pairing of a distribution with
the spectral transform, not as an integral with a scalar function $m$.
\end{example}

The paper is organized as follows.  Section~\ref{sec:lit} distinguishes
sources available by 1994 from later developments.  Background and
notation are given in Sections~\ref{sec:background} and
\ref{sec:notations}.  Sections~\ref{sec:orbital} and
\ref{sec:unitary} review orbital integrals and rank-one representations.
Section~\ref{sec:invariant} separates algebraic trace density from the
conditional topological spectral statement.  We end with a summary,
open issues, and a glossary.

\section{Literature review}\label{sec:lit}
The following sources are grouped by their relation to the original
1994 manuscript rather than presented as a claim of an exhaustive survey.
\paragraph{Through 1994.}
Kubota's construction and the theory of central extensions underlie the
rank-one cover (see the sources assembled in \cite{Jo94}).  Gelbart's
account \cite{Ge76}, Kazhdan--Patterson's metaplectic forms \cite{KP84},
and Flicker--Kazhdan's correspondence \cite{FK86} provide much of the
representation-theoretic setting.  Arit\"urk studied a threefold cover
\cite{Ar80}, and Moen analyzed irreducibility of unitary principal series
\cite{Mo88}.  Harish-Chandra's harmonic analysis \cite{HC70,HC78},
Kazhdan's density ideas \cite{Ka86}, and Arthur's trace Paley--Wiener
work \cite{Ar94} motivate the invariant-distribution question.  These
results concern different groups and test-function spaces; citations
alone are not substitutes for the hypotheses needed below.
\paragraph{After 1994.}
The 2001 paper \cite{Jo01} supplied the immediate source for many of the
rank-one formulations below.  The 2012 work of Wen-Wei Li
\cite{Li12}, especially Sections~5.4, 5.6, and Theorem~5.8.10,
places tempered Fourier analysis and density for covering groups in the
setting of a local invariant trace formula.  Li states a non-Archimedean
$R$-group decomposition as an admitted hypothesis in Theorem~5.4.1;
its role must be kept visible when invoking the later results.

\begin{example}[A chronological distinction]
Moen's 1988 irreducibility results \cite{Mo88} could inform the 1994
notes.  Li's 2012 tempered-character density theorem \cite{Li12}
could not, and belongs to the modern reassessment rather than to
the original claims.
\end{example}

\section{Background}\label{sec:background}
We fix the cover and the convention for its central character before
describing invariant harmonic analysis.  Except where an elementary
example is explicitly noted, subsequent statements concern the genuinely
$\mu_n$-equivariant component.


\subsection{Assumptions on \texorpdfstring{$\overline{G}$}{the covering group}}\label{section:assume}

Let $F=\F$ be a finite extension of $\mathbb Q_p$ with residual
characteristic $p$.  Fix an integer $n>1$ satisfying $p\nmid 2n$
and assume $\mu_n\subset F^\times$.  Set $G=SL(2,F)$.  The Kubota
$n$-fold metaplectic cover is a finite topological central extension
\[
1\longrightarrow\mu_n\longrightarrow\bG\xrightarrow{\ \pr\ }G
\longrightarrow1.
\]
The restriction $p\nmid 2n$ is stronger than the original 1994
manuscript's generality; no assertion for wildly ramified covers is
made here.


Let $\varpi$ be a uniformizer, $\mathcal O$ the ring of integers, $q=|\mathcal O/\varpi\mathcal O|$,
and $|\cdot|$ the normalized absolute value. Put
\begin{equation}\label{eqn:N}
\nu=\begin{cases}n,& n\text{ odd},\\ n/2,& n\text{ even},\end{cases}
\end{equation}
and let $N_0$ denote
the upper triangular unipotent subgroup of $G$.

For $g=\begin{psmallmatrix}a&b\\c&d\end{psmallmatrix}\in SL(2,\F)$ let $x(g)=c$ if $c\ne0$ and
$x(g)=d$ if $c=0$, and let $\beta:G\times G\to\mu_n$,
\[
\beta(g_1,g_2)=(x(g_1),x(g_2))_n\,(-x(g_1)^{-1}x(g_2),x(g_1g_2))_n
=\Bigl(\tfrac{x(g_1g_2)}{x(g_1)},\tfrac{x(g_1g_2)}{x(g_2)}\Bigr)_n,
\]
where $(\cdot,\cdot)_n$ denotes the $n$th power Hilbert symbol
\cite{We75}.  The extension has the group law
\begin{equation}\label{eq:group-law}
 (g_1,\zeta_1)(g_2,\zeta_2)=(g_1g_2,\beta(g_1,g_2)\zeta_1\zeta_2).
\end{equation}
This is the local cocycle model used in the 1994 manuscript
\cite{Jo94}.
\begin{example}[The central kernel]\label{ex:central}
For $\zeta\in\mu_n$, the element $(1,\zeta)$ is central.  If $\pi$
is genuine then $\pi(1,\zeta)=\zeta\,\mathrm{id}$, while a genuine
test function satisfies $f((1,\zeta)x)=\zeta^{-1}f(x)$.
Thus $f(x)\pi(x)$ is unchanged under changing the lift of a point
of $G$.
\end{example}

\subsection{Basic notation}\label{subsec:notation}

For $H\subset G$ put $\overline H=\pr^{-1}(H)$; in particular $\overline{G_r}$ is the preimage of the set $G_r$ of regular semisimple elements of $SL(2,F)$. Let $\mathcal L(G)=\{G,A\}$ denote the standard Levi subgroups, where $A$ is the diagonal torus, and write $\mathcal L(\bG)=\{\bG,\bA\}$ for their preimages. For a Levi
$M$ let $X(M)$ (resp.\ $X^{un}(M)$) be the variety of unramified (resp.\ unramified unitary)
characters; for $M=\bA$ these are the characters of $A^n$ (identified with characters of
$\overline{A^n}$). Let $W=N_G(A)/A$ be the Weyl group. For $G=SL(2,\F)$,
\[
w_0=\begin{pmatrix}0&-1\\1&0\end{pmatrix}\in G
\]
represents the nontrivial element of the order-two group $W$.  Although $w_0^2=-I\in A$, the representative $w_0$ itself need not have order two. The
group $W$ acts on $\bA$ by conjugation by any lift of $w_0$; we use the lift $\bw_0=(w_0,1)$.

The unipotent radical $N$ of a parabolic subgroup is identified with $\{(u,1)\mid u\in N\}\subset\bG$.
Let $K_0=SL(2,\mathcal O)$. For $SL(2)$ and $p\nmid 2n$, the cover splits over $K_0$ \cite{Ge76}.

\paragraph{Genuine functions and representations.} A function $f$ on $\bG$ (or on a subgroup
$H\subset \bG$) is \emph{genuine} if $f(g,\varsigma)=\varsigma^{-1}f(g,1)$. A representation $\pi$ is
\emph{genuine} if $\pi(g,\varsigma)=\varsigma\,\pi(g,1)$. For genuine $f$ and $\pi$ the product
$f(x)\pi(x)$ is invariant under $\mu_n$, so $\pi(f)=\int_G f(\tilde g)\pi(\tilde g)\,dg$ is defined
(for any section $g\mapsto\tilde g$). Let $C_c^\infty(\bG)$ be the space of smooth (locally constant,
compactly supported) genuine functions. Write $f^y(x)=f(y^{-1}xy)$ and
\[
\sJ=\Span\{f-f^y\mid f\in C_c^\infty(\bG),\ y\in\bG\}.
\]
A linear functional $D$ is \emph{invariant} if $D(f^y)=D(f)$ for all $y$, i.e.\ $D|_{\sJ}=0$.

\paragraph{The Schwartz space (corrected definition).} Let $\|g\|=\max|g_{ij}|$ and put
\[\sigma(g)=\log\max\bigl(\|g\|,\,\|g^{-1}\|\bigr)\ \ge0.\] Fix the maximal compact subgroup $K_0$ and the standard Borel
subgroup $B$. Let $\delta_B$ be the modulus character of $B$, extended to $G=BK_0$
(Iwasawa decomposition) by $\delta_B(bk)=\delta_B(b)$. Harish-Chandra's function is
\[
\Xi(x)=\int_{K_0}\delta_B(xk)^{-1/2}\,dk,
\]
pulled back to $\bG$ via $\pr$. For a compact open subgroup $K\subset\bG$ with $K\cap\mu_n=1$
(so that bi-$K$-invariant genuine functions exist) let $\calC_K(\bG)$ be the space of genuine
functions $f$ on $\bG$ that are bi-$K$-invariant (hence locally constant) with
\[
\upsilon_k(f)=\sup_{x\in\bG}|f(x)|\,\frac{(1+\sigma(x))^k}{\Xi(x)}<\infty\qquad(k\ge0).
\]
Topologize $\calC_K(\bG)$ by the seminorms $\upsilon_k$ and put
$\calC(\bG)=\bigcup_K\calC_K(\bG)$ with the locally convex inductive limit topology. This is the
\emph{Harish-Chandra Schwartz space} (in the $p$-adic case the letter $\calC$ is often used for the Schwartz space, rather than $\cS$;
compact support is not part of the definition). It is an algebra under convolution
and $C_c^\infty(\bG)\subset\calC(\bG)\subset L^2(\bG)$ with $C_c^\infty$ dense.
A \emph{tempered distribution} on $\bG$ is a continuous linear functional on $\calC(\bG)$.

\paragraph{Characters and dual objects.} For an admissible genuine representation $\pi$ of $\bG$ let
$\Theta_\pi$ be its character, regarded as a locally integrable genuine function on $\overline{G_r}$ or
as an invariant distribution (see \cite{Jo94}, \S4, for the extension of \cite{Si79}, \S4.8, to covers;
\cite{HC70,HC78}). For a Levi $\overline M$ let $\Pi(\overline M)$ be the genuine admissible dual,
$\Pi(\overline M)_t$ the genuine tempered dual, $\Pi(\overline M)_u$ the genuine unitary dual
(equivalence classes of irreducible representations; all representations are admissible).
Write $\Pit=\Pi(\bG)_t$. The \emph{trace transform} is
\[
\sF:C_c^\infty(\bG)\to\C^{\Pit},\qquad \sF f(\pi)=\Theta_\pi(f)=\tr\pi(f).
\]
Let $R_t(\bG)$ be the Grothendieck group of genuine, tempered, admissible representations of $\bG$,
tensored with $\C$. The image of $\sF$, viewed as a space of functionals on $R_t(\bG)$, is the
\emph{tempered Paley--Wiener space} $PW_t(\bG)$. Similarly $PW(\bG)$ denotes the (non-tempered)
Paley--Wiener space, the space of functionals $\pi\mapsto\Theta_\pi(f)$ on $R(\bG)$, $f\in C_c^\infty(\bG)$.

%====================================================================

\section{Notation and conventions}\label{sec:notations}\label{subsec:notsum}

For convenience, the main notation is collected here.  Every symbol also receives an explanation at its first substantive use.

{\small
\begin{longtable}{@{}>{\raggedright}p{4.1cm}>{\raggedright\arraybackslash}p{11.2cm}@{}}
\toprule
\textbf{Symbol} & \textbf{Meaning} \\
\midrule
\endhead
$\F$, $\mathcal O$, $\varpi$, $q$, $|\cdot|$, $p$ & $p$-adic field; its ring of integers; a uniformizer; the order of the residue field; the normalized absolute value; the residual characteristic \\
$n$, $\mu_n$ & degree of the cover; the group of $n$-th roots of unity (contained in $\F$) \\
$G$, $\bG$, $\pr$ & $SL(2,F)$, its $n$-fold cover, and the projection $\bG\to G$ \\
$\nu$ & $n$ if $n$ is odd, $n/2$ if $n$ is even (equation \eqref{eqn:N}) \\
$A$, $\bA$, $N_0$ & diagonal subgroup of $G$; its preimage in $\bG$; the upper triangular unipotent subgroup \\
$W$, $w_0$, $\bw_0$ & Weyl group $N_G(A)/A$; the representative $\begin{psmallmatrix}0&-1\\1&0\end{psmallmatrix}$ for $SL(2)$; its lift $(w_0,1)$ \\
$(\cdot,\cdot)_n$, $\beta$ & Hilbert symbol; the cocycle defining $\bG$ \\
$K_0$, $K$ & the maximal compact subgroup $SL(2,\mathcal O)$; compact open subgroups of $\bG$ meeting $\mu_n$ trivially \\
$\|g\|$, $\sigma(g)$ & max-norm of the matrix entries; the height $\log\max(\|g\|,\|g^{-1}\|)$ \\
$\delta_B$, $\Xi$ & modulus character of the Borel subgroup $B$; Harish-Chandra's function \\
$C_c^\infty(\bG)$ & smooth, compactly supported, genuine functions on $\bG$ \\
$\calC_K(\bG)$, $\calC(\bG)$ & $K$-bi-invariant genuine functions with $\Xi$-weighted rapid decay; the Schwartz space $\bigcup_K\calC_K(\bG)$ \\
$f^y$, $\sJ$ & $f^y(x)=f(y^{-1}xy)$; the span of all $f-f^y$ in $C_c^\infty(\bG)$ \\
$\Theta_\pi$ & the character of an admissible genuine representation $\pi$ \\
$\Pi(\overline M)$, $\Pi(\overline M)_t$, $\Pi(\overline M)_u$ & genuine admissible, tempered, unitary dual of the Levi subgroup $\overline M$; $\Pit=\Pi(\bG)_t$ \\
$R_t(\bG)$, $PW_t(\bG)$, $PW(\bG)$ & Grothendieck group of genuine tempered representations (tensored with $\C$); the tempered Paley--Wiener space; the Paley--Wiener space \\
$\sF$, $\widehat D$ & the trace transform $f\mapsto(\pi\mapsto\Theta_\pi(f))$; the Fourier transform of an invariant distribution $D$, defined by $D=\widehat D\circ\sF$ \\
$I_M(\tau)$, $I_{\bA}(\chi)$, $V_\tau$ & unitarily induced representations; the space of the supercuspidal $\tau$ \\
$\chi^{\bw}$ & the conjugate $\chi^{\bw}(a)=\chi(\bw^{-1}a\bw)$ \\
$C$, $\bC$, $\mu$, $\chi_\mu$, $s(\mu)$ & maximal subgroup of $A$ with $\bC$ abelian; a character of $\bC$; $\Ind_{\bC}^{\bA}\mu$; the exponent in $\mu(x)=\mu_0(x)|x|^s$ \\
$K(\mu)$, $V(\mu)$, $T_w$, $J_\mu$, $d$, $\mathcal R$ & spaces and intertwining operator of \S\ref{subsec:cs}; the matrix of the operator $J$; the cardinality and a set of representatives of $\bA/\bC$ \\
$\pi_{sp}$, $\pi_{nt}$ & the special representation; the non-tempered (``trash'') quotient \\
$D(t)$, $Z_t$, $F_f^{Z_t}$, $F_f^{A^\nu}$ & Weyl discriminant; centralizer of $t$ in $\bG$; orbital integrals (\S\ref{sec:orbital}) \\
$\Delta$ & Laplacian on a compact torus (\S\ref{subsec:main}) \\
\bottomrule
\end{longtable}}


\begin{example}[A Weyl representative]
For $a\in F^\times$, conjugation by the matrix $w_0$ sends
$\operatorname{diag}(a,a^{-1})$ to
$\operatorname{diag}(a^{-1},a)$.  Hence a Weyl-invariant test
function on the diagonal torus satisfies $h(a)=h(a^{-1})$.
\end{example}

\section{Basic lemmas on orbital integrals}\label{sec:orbital}

If $p<\infty$ then $\F^\times=\varpi^{\Z}\cdot\mu_{q-1}\cdot U_1$ (a direct product). If
$p\nmid 2n$ then $\mu_n\subset\F^\times$ gives $q\equiv1\pmod n$. Recall $\nu$ from \eqref{eqn:N}.

\begin{lemma}\label{lem:subgroup}
Let $A$ be the diagonal subgroup of $SL(2,F)$ and let
$h(a)=\operatorname{diag}(a,a^{-1})$.  Under the standing hypotheses,
\[
 C=\varpi^{\mathbb Z}\mathcal O^{\times\nu}
 =\varpi^{\mathbb Z}\mu_{q-1}^{\nu}(1+\varpi\mathcal O)
\]
is a subgroup of index $\nu$ in $F^\times$.
\end{lemma}
\noindent\textbf{Proof.}
The unit group decomposes as
$\mathcal O^\times=\mu_{q-1}\times(1+\varpi\mathcal O)$.
Because $p\nmid\nu$, raising to the $\nu$th power is an automorphism
of $1+\varpi\mathcal O$.  Also $n\mid(q-1)$ and $\nu\mid n$, so
$\nu\mid(q-1)$.  Since $\mu_{q-1}$ is cyclic, its subgroup of
$\nu$th powers has index $\nu$.  The uniformizer factor contributes
no additional index.\hfill$\square$

\begin{remark}[Maximal abelian preimages]\label{rem:maxabelian}
The earlier papers additionally identify $C$ as a maximal subgroup
of $F^\times$ with abelian preimage $\overline{h(C)}$ in the covering
diagonal torus \cite{Jo94,Jo01}.  This stronger assertion depends on
the Hilbert-symbol commutator for the chosen cocycle and has not been
reproved here.  It is not needed for Lemma~\ref{lem:subgroup}.
\end{remark}

Let $D(h(a))=\det(1-\mathrm{Ad}(h(a)))_{\mathfrak g/\mathfrak a}=-(a-a^{-1})^2$, which pulls back to
$\overline{G_r}$ via $\pr$. For $t\in\overline{G_r}$ let $Z_t=\Cent(t,\bG)$. (This need not be the cover of
a centralizer in $G$: for $t=(x,1)$ in general $\Cent(t,\bG)\ne\overline{\Cent(x,G)}$.) The
\emph{orbital integral} of $f\in C_c^\infty(\bG)$ is
\begin{equation}\label{eqn:orb}
F_f^{Z_t}(t)=|D(t)|^{1/2}\int_{Z_t\backslash\bG}f(x^{-1}tx)\,\frac{dx}{dt}
\end{equation}
(it exists by Harish-Chandra \cite{HC70}; $D$ is defined by identifying $Z_t$ with a torus over
$\overline\F$). For regular $a\in\bA$ put
$F_f^{A^\nu}(a)=|D(a)|^{1/2}\int_{\overline{A^\nu}\backslash\bG}f(x^{-1}ax)\,\frac{dx}{da}$.

\begin{remark}[Orbital-integral statements from the earlier sources; verification pending]\label{lem:orbital}
Let $\bG$ be a cover of $SL(2,\F)$ as in \S\ref{section:assume}.
\begin{enumerate}[label=(\alph*),leftmargin=*,itemsep=1pt]
\item For $f\in C_c^\infty(\bG)$ and $a\in\bA-\overline{A^\nu}$, $F_f^{A^\nu}(a)=0$.
\item $f\mapsto F_f^{A^\nu}$ is a surjection $C_c^\infty(\bG)\to C_c^\infty(\overline{A^\nu})^W$, where $W$
acts on $\overline{A^\nu}$ as in \S\ref{subsec:notation}.
\item $f\mapsto F_f^{A^\nu}$ is a continuous surjection $\calC(\bG)\to\calC(\overline{A^\nu})^W$.
\end{enumerate}
\end{remark}
These assertions reproduce the earlier formulation.  The hypotheses and
the claimed surjectivity for the genuine Schwartz space have not been
checked against the original proofs in this revision, so they are recorded
as source-dependent statements rather than newly proved lemmas.  See
\cite{Jo94,Jo01}. 

%====================================================================
\begin{example}[The discriminant]
For $h(a)=\operatorname{diag}(a,a^{-1})$ with $a^2\ne1$,
the two nontrivial adjoint eigenvalues are $a^2$ and $a^{-2}$.
Thus the rank-one Weyl discriminant is
$(1-a^2)(1-a^{-2})=-(a-a^{-1})^2$, as used in
\eqref{eqn:orb}.  The restriction $a^2\ne1$ is exactly the
regularity condition on the split torus.
\end{example}

\section{Unitary and tempered dual of \texorpdfstring{$\overline{SL(2,\F)}$}{SL(2,F) cover}}\label{sec:unitary}

In this section $\bG$ is a cover of $SL(2,\F)$ as in \S\ref{section:assume}. The results are
essentially in the literature (see \cite{Jo99a}, \cite{Jo01}); they are recalled in order to describe
the tempered dual in \S\ref{sec:invariant}. 

It is remarked in \cite{Be84}, \S2.2 that the arguments of \cite{BZ76}, Ch.~2 carry over to finite
central extensions of reductive groups over $p$-adic fields (see also \cite{KP84}, \S1.2), and the
arguments of \cite{BZ77}, \S2 carry over to split reductive groups. These reduce the determination of the
unitary dual to classifying the supercuspidal representations (done in \cite{Jo99b} for
$\gcd(p,n)=1$ and in \cite{Jo98} for any $p,n$) and the constituents of induced representations.

\subsection{Principal series}\label{subsec:ps}

Let $\overline P=\overline MN$ be a Levi decomposition and $(\tau,V_\tau)$ a genuine
supercuspidal representation of $\overline M$ (extended trivially to $N$). The \emph{unitarily
induced representation} $I_M(\tau)$ is the right-translation representation on
\begin{multline*}
\Bigl\{f:\bG\to V_\tau\ \text{genuine}\ \Big|\ f(mg)=\delta_P(m)^{1/2}\tau(m)f(g),\\
f\text{ right-invariant under some compact open }K\Bigr\}.
\end{multline*}
If $\varrho$ is a representation of $H\subset\bG$ and $x\in N_{\bG}(H)$, put
$\varrho^x(h)=\varrho(x^{-1}hx)$.

For $\chi,\chi'\in\Pi(\bA)$, call $\chi$ \emph{regular} if $\chi^{\bw}\ne\chi$ for all
$w\in W-\{1\}$, and call $\chi,\chi'$ \emph{$W$-conjugate} if $\chi'=\chi^{\bw}$ for some $w$.
Distinct $W$-conjugacy classes yield inequivalent representations. A \emph{unitary principal series} representation is an induced representation
$I_{\bA}(\chi)$ with $\chi\in\Pi(\bA)_u$. It is tempered, but need
not be irreducible; only its irreducible constituents determine points
of $\Pi(\bG)_t$.  When $I_{\bA}(\chi)$ is reducible, we call those
constituents \emph{reducible principal-series constituents}.

\begin{proposition}[Moen \cite{Mo88}]\label{prop:moen}
Assume $\gcd(p,n)=1$.
\begin{enumerate}[label=(\alph*),leftmargin=*,itemsep=1pt]
\item If $n$ is even, $I_{\bA}(\chi)$ is irreducible and unitary for every $\chi\in\Pi(\bA)_u$.
\item If $n$ is odd, $I_{\bA}(\chi)$ is irreducible and unitary for every $\chi\in\Pi(\bA)_u$
with $\chi^{\bw_0}\ne\chi$. If $\chi^{\bw_0}=\chi$ then $I_{\bA}(\chi)$ is reducible with two irreducible
constituents, except possibly for one exceptional $\chi$ singled out in \cite{Mo88}.
\end{enumerate}
\end{proposition}
%\verify{The earlier version states the exception as ``$\chi=1$''. For a genuine character of $\bA$ the
%trivial character is not an option, so the exceptional datum must be re-expressed (for example in
%terms of the associated character of $C$ or of $A^n$). Check against \cite{Mo88}.}
Moen \cite{MoND} computes the intertwining operators explicitly, as matrices in the Kirillov model, for
$\gcd(p,n)=1$. The irreducibility statement has a direct Bruhat-theory proof and can also be deduced from
\cite{FK86}.

\subsection{Complementary series}\label{subsec:cs}

We review results of Arit\"urk \cite{Ar80} and use \cite{FK86} to extend them to the $n$-fold cover
(\cite{Ar80} assumed $n=3$ and $p>3$; for $n=2$ the results were essentially known to Gelbart--Sally
\cite{GS75}). An irreducible unitary representation is called a \emph{complementary series} representation if it occurs as
$I_{\bA}(\chi)$ with $\chi\in\Pi(\bA)-\Pi(\bA)_u$; these are not tempered.

Let $C\subset A$ be a maximal subgroup with $\bC\subset\bA$ abelian. For $\mu\in\Pi(\bC)$ let
$\chi_\mu=\Ind_{\bC}^{\bA}\mu$. If $\mu(x)=\mu_0(x)|x|^s$ with $\mu_0$ of finite order and
$s\in\C$, write $s=s(\mu)$. Let $K(\mu)$ be the space of locally constant $f:\F\times\bA\to\C$ with
(i) $f(x,a_1a_2)=\mu(a_1)f(x,a_2)$ for $a_1\in\bC$, $a_2\in\bA$; and
(ii) $|x|\,\chi(h(x),1)f(x,a)$ constant for $|x|$ large. Let $\mathcal R\subset\bA$ be a set of
representatives for $\bA/\bC$, of cardinality $d$ ; then $f\in K(\mu)$ is
identified with the $d$-tuple $(f(x,a))_{a\in\mathcal R}$.

Let $V(\mu)$ be the space of locally constant $\vphi:\bG\times\bA\to\C$ with
(i) $\vphi(g,a_1a_2)=\mu(a_1)\vphi(g,a_2)$, $a_1\in\bC$; and
(ii) $\vphi(a_1ug,a_2)=\delta(a_1)\vphi(g,a_2a_1)$, $a_1,a_2\in\bA$, $u\in N_0$.
For $w\in W$ define the intertwining operator
\[
T_w\vphi(g,a)=\int_{\F}\vphi\bigl(\bw\cdot(\begin{psmallmatrix}1&x\\0&1\end{psmallmatrix},1)\cdot g,\ \bw a\bw^{-1}\bigr)\,dx,
\qquad \mathrm{Re}\,s(\mu)>0.
\]
\begin{lemma}[Arit\"urk]\label{lem:T} $T_w$ intertwines $I_{\bA}(\mu)$ and $I_{\bA}(\mu^{\bw})$.\end{lemma}
\begin{lemma}[Arit\"urk]\label{lem:contra} The pairing
$\langle\vphi_1,\vphi_2\rangle=\int_{\bar B\backslash\bG}\int_{\bA/\bC}\vphi_1(g,a)\vphi_2(g,a)\,da\,dg$
is a non-degenerate bilinear form on $V(\mu)\times V(\mu^{\bw})$; hence $I_{\bA}(\mu^{\bw})$ is the
contragredient of $I_{\bA}(\mu)$.\end{lemma}
(Here $\bar B$ is the Borel subgroup of $\bG$; the integrand lies in the space $L(\bG,\bar B)$ of functions
with $\vphi(\begin{psmallmatrix}a&*\\0&a^{-1}\end{psmallmatrix},\varsigma)g)=|a|^2\vphi(g)$.) Neither
lemma needs $\gcd(p,n)=1$.

For $f\in K(\mu)$ define $f^\wedge(x,a)=\int_{\F^\times}f(y,a)\psi(xy)\,dy$ with $\psi$ a fixed additive
character.
\begin{lemma}[Arit\"urk]\label{lemma:Jdef}
Assume $\gcd(p,n)=1$, $\mu(x)=|x|^s$. For $\vphi_1,\vphi_2\in V(\mu)$,
\[
\langle\vphi_1,T\vphi_2\rangle=\int_{\F}\int_{\bA/\bC}f_1^\wedge(x,a)\,\overline{(Jf_2^\wedge)(x,a)}\,da\,dx,
\]
where $J=J_\mu$ is a linear transformation of $K(\mu)^\wedge$ (an $d\times d$ matrix) and
$f_i(x,a)=\vphi_i(\bw^{-1}\cdot(\begin{psmallmatrix}1&x\\0&1\end{psmallmatrix},1),a)$.
\end{lemma}
%\verify{As written, the left side is complex-bilinear while the right side is conjugate-linear in
%$\vphi_2$; the left side should be a sesquilinear pairing obtained from Lemma~\ref{lem:contra} by the
%identification of $\overline{V(\mu)}$ with $V(\bar\mu)$. Consult \cite{Ar80} for the exact normalization.}

\begin{lemma}[Langlands, Arit\"urk] Assume $\gcd(p,2n)=1$. If $0\le\mathrm{Re}\,s(\mu)\le1/n$ and
$|\mathrm{Im}\,s(\mu)|\le\pi/(n\ln q)$ then the image of $J_\mu$ is an irreducible representation of $\bG$.\end{lemma}
\begin{proposition}[Flicker--Kazhdan \cite{FK86}]\label{prop:fk}
If $0<s(\mu)<1/n$ then $I_{\bA}(\mu)$ is unitarizable.\end{proposition}
\begin{corollary}
If $0<s<1/n$ and $\mu(x)=|x|^s$, then $I_{\bA}(\mu)$ is unitary (unitarizable) in this range; the Hermitian
form obtained from the intertwining pairing is positive definite, and if in addition $\gcd(p,2n)=1$ the
matrix $J_\mu$ is positive definite.
%\verify{the earlier version called $\langle\vphi_1,T\vphi_2\rangle$
%itself ``positive definite''; it is bilinear (indeed $\langle iv,T\,iv\rangle=-\langle v,Tv\rangle$ for a
%bilinear form), so only the Hermitian form can be positive definite. Check the claim about $J_\mu$ against
%\cite{Ar80}}
\end{corollary}

\subsection{The special and the ``trash'' representations}\label{subsec:special}

For the origin of the term ``trash'' representation, see \cite{Ge76}.

\begin{remark}[Endpoint of the complementary series: an unresolved source conflict]\label{prop:special}
At $s=1/n$ with $\mu(x)=|x|^s$, the 1994 manuscript
\cite{Jo94} conjectured that the \emph{image} of the intertwining
matrix $J_\mu$ was the square-integrable special representation
$\pi_{sp}$, while its \emph{kernel} was the non-tempered constituent
$\pi_{nt}$.  The supplied corrected 2026d draft asserted the opposite:
$\ker J_\mu\simeq\pi_{sp}$ and
$\operatorname{im}J_\mu\simeq\pi_{nt}$.  These opposite identifications
cannot both be retained as theorems without specifying and comparing
the normalization, domain, and direction of $J_\mu$.  The notation
$\pi_{sp},\pi_{nt}$ below refers to the respective square-integrable
and non-tempered constituents without identifying which is the
kernel or image of this operator.

The 2026d draft cites Arit\"urk \cite{Ar80}, Gelbart--Sally
\cite{GS75}, and Kazhdan--Patterson \cite{KP84} concerning endpoint
constituents and unitarity.  Those citations have not been checked here
for the exact general-$n$ assertions and the possible interchange
between subrepresentation and quotient.  The endpoint identification
must therefore be resolved against the original intertwining-operator
conventions before claiming a corrected classification.
\end{remark}

\subsection{Classification}

\begin{theorem}\label{thm:class-u}
Let $\bG$ be as in \S\ref{section:assume} (cover of $SL(2)$). If $\pi\in\Pi(\bG)_u$ then one of the following
holds.
\begin{itemize}[leftmargin=*,itemsep=1pt]
\item (Principal series) $\pi=I_{\bA}(\chi)$ with $\chi\in\Pi(\bA)_u$, $\chi^{\bw_0}\ne\chi$ (\S\ref{subsec:ps}).
\item (Complementary series) $\pi=I_{\bA}(\chi)$ with $\chi\in\Pi(\bA)-\Pi(\bA)_u$ (\S\ref{subsec:cs}).
\item (Reducible principal series; occurs only for $n$ odd by Proposition~\ref{prop:moen}) $\chi\in\Pi(\bA)_u$ with
$\chi^{\bw_0}=\chi$, and $\pi$ is a subrepresentation or quotient of $I_{\bA}(\chi)$.
\item $\pi$ is a ``special'' or ``trash'' representation (\S\ref{subsec:special}).
\item $\pi$ is supercuspidal (\cite{Jo99b}).
\end{itemize}
\end{theorem}
\begin{theorem}\label{thm:class-t}
If $\pi\in\Pi(\bG)_t$ then $\pi$ is one of: a principal series $I_{\bA}(\chi)$ with $\chi\in\Pi(\bA)_u$,
$\chi^{\bw_0}\ne\chi$; a constituent of a reducible $I_{\bA}(\chi)$ with $\chi$ unitary and $\chi^{\bw_0}=\chi$;
the special representation $\pi_{sp}$; or a supercuspidal representation. In particular
$\Pi(\bG)_t\subset\Pi(\bG)_u$, and for $n$ even there are no reducible unitary principal series.
\end{theorem}
For the inclusion $\Pi(\bG)_t\subset\Pi(\bG)_u$ in this rank-one setting, see the case-by-case discussion in \cite{Jo94,Jo01}.  The precise completeness of the list is not independently verified here.

\begin{remark}[Classification status in this edition]\label{rem:classification-status}
The preceding classification is retained as a guide to the expected
rank-one structure and to the assertions in \cite{Jo94,Jo01}.  The
exact exceptional parameters in Proposition~\ref{prop:moen}, the
normalization of the intertwining Hermitian form, and the special and
non-tempered endpoint constituents have not been rechecked against
all primary sources in this revision.  The statements are therefore
not presented as new, independently verified classification theorems.
\end{remark}

\begin{example}[An admissible tame datum]\label{ex:tame}
Take $F=\mathbb Q_7$ and $n=3$.  Since $3\mid(7-1)$, the field
contains $\mu_3$, and $7\nmid 2n$.  The invariant $\nu$ of
\eqref{eqn:N} is $3$.  This provides a concrete rank-one setting
in which the covering torus and its genuine inducing characters
can be studied using the unitary circles in the unramified-character
parameter space.
\end{example}

\section{The 1994 conjectures and the present scope}\label{sec:conjectures}
The 1994 paper is more than an abbreviated version of the 2001 article.
It develops the local covering torus, Jacquet functors, characters,
and several spectral questions under explicitly conjectural hypotheses.
This section records their disposition in the present narrower treatment.

\begin{remark}[Principal-series reducibility in the 1994 draft]
The conjectures in the 1994 section ``Principal series, reducible
principal series'' include irreducibility for even $n$ and a reducibility
pattern for odd $n$ \cite{Jo94}.  The corrected 2026d paper replaces
these by an attribution to Moen \cite{Mo88}, but retains an exceptional
parameter whose precise genuine-character meaning needs verification.
This edition does not treat the 1994 conjectural wording as a proof.
\end{remark}

\begin{remark}[Intertwining operators and endpoint conventions]
The original conjectures concern the determinant, positivity, and
kernel and image of the matrix $J_\mu$ attached to a rank-one
intertwining operator.  The strongest caution is the reversed
kernel--image identification documented in Remark~\ref{prop:special}.
Even if both manuscripts describe the same composition factors,
these conventions must be reconciled before asserting the result.
\end{remark}

\begin{remark}[Temperedness and Paley--Wiener assertions]
The 1994 typescript labels tempered unitarity and some Fourier-image
theorems as conjectural or as dependent on the conjectural dual
classification.  The later density theorem of
\cite[Theorem~5.8.10]{Li12} addresses an algebraic trace-kernel
question.  It does not, by itself, establish the particular
principal-series surjectivity formulas or the full topological
claim of Theorem~\ref{thm:main}.
\end{remark}

\begin{example}[Revising a conjectural assertion]\label{ex:conjecture-status}
Let $\chi$ be a unitary genuine inducing datum fixed by the Weyl
action.  A relation for the trace of $I_{\bA}(\chi)$ only gives
the sum of the traces of its irreducible constituents.  A correct
Fourier-image theorem for irreducible tempered characters must also
specify compatibility with the individual constituents at this point.
The conjectural Fourier-image formulas of the original guide do not
address that distinction explicitly.
\end{example}

%====================================================================

\section{Invariant distributions}
\label{sec:invariant}

We classify the image of $C_c^\infty(\overline{G})$
and of ${\calC}(\overline{G})$ under the ``scalar-valued
Fourier transform'' or ``trace map'',
$f\longmapsto {\rm tr}\, \pi(f)$, where $\overline{G}$ is as
in \S \ref{section:assume} above. 
We describe the algebraic factorization through tempered traces, where $\overline{G}$ is a cover of $SL(2,{\F})$ as in Section~\ref{section:assume}.
Finally, under an explicitly identified topological Paley--Wiener input,
we obtain a continuous spectral functional for invariant tempered
distributions.  This need not be an integral against a scalar density.

\subsection{Tempered Paley-Wiener theorem}
\label{subsec:TPW}

In this section, $\overline{G}$ is a cover of $SL(2,{\F})$ as in \S \ref{section:assume}.
Let $C\subset A$ be a maximal subgroup of $A$ for which 
$\overline{C}\subset \overline{A}$ is abelian.
Next we classify the image
of the Fourier transforms of a ``generic'' {\em unitary} principal series
representation 
$
\pi =I_{\overline{A }}(\chi ),
$
(where $\chi =Ind_{\overline{C}}^{\overline{A }}\mu \in \Pi (\overline{A }),
\mu \in \Pi (\overline{C })$),
on $C_c^\infty (\overline{G })$.
Both the Weyl group $W$ and
$H=\overline{A}/\overline{C}$ act on functions on $\Pi(\overline{C})$ by conjugation.
Let $C^\infty (\Pi(\overline{C})_u)^{WH}$ denote the 
$W$- and $H$-invariant smooth functions on the compact torus $\Pi(\overline{C})_u$.

Recall the Fourier transform with respect to the 
principal series,
$
\phi_f(\mu)=\Theta _\pi(f),
$
(where $\pi =I_{\overline{A}}(\chi ),
 \chi =Ind_{\overline{C}}^{A }\mu
\in \Pi(\overline{A }),
\mu\in \Pi(\overline{C })$).
When $\mu$ is unitary we call this the
Fourier transform with respect to the 
{\it unitary principal series}.
If the restriction of $\mu$ to an
$diag(x,x^{-1})\in A^n$ is of the form $|x|^s$ then
we write $\phi_f(\mu)=\phi_f(s)$.
When $\mu$ is of this form and $s$ is real, 
we call this the
Fourier transform with respect to the 
{\it complementary series}.

\begin{remark}[Principal-series images; verification pending]\label{prop:image}
For $f\in C_c^\infty (\overline{G })$, the image
$C_c^\infty (\overline{G })_{ps}^{\wedge }$
of the Fourier transform $f\longmapsto \phi_f$
with respect to the unitary principal series,
is given by 
\[
C_c^\infty (\overline{G })_{ps}^{\wedge }
=\left\{h\in C^\infty(\Pi (\overline{C})_u)^{WH}\;|\;
\begin{array}{c}
h\;{\rm is\;a\;trig\;polynomial\;on} \\ 
{\rm each\;circle\;in\;}\Pi(\overline{C})_u
\end{array}
\;\right\}.
\]
The image
$C_c^\infty (\overline{G })_{cs}^{\wedge }$
of the Fourier transform $f\longmapsto \phi_f$
with respect to the complementary series,
is given by 
\[
C_c^\infty (\overline{G })_{cs}^{\wedge }
\cong \left\{h\ {\rm restricted \ to}\ 0<s<1/n\ \Big|\ h\ {\rm a}\ WH\mbox{-invariant\ Laurent\ polynomial\ in\ }q^{\pm s}\right\}.
\]
\end{remark}

These formulas are reproduced from the earlier paper as the
\emph{proposed} principal-series images, not as independently verified
surjectivity results.  In particular, finite support across the
disconnected Bernstein components and compatibility at reducibility
parameters require separate checks.  The motivating argument uses character formulas for induced representations and results on $p$-adic Mellin transforms
in \cite{Ta75}, pp. 43-44. Further details in the earlier formulation appear in \cite{Jo01}. (Since $W$ acts by $s\mapsto-s$, a $W$-invariant Laurent
polynomial in $q^{\pm s}$ is a polynomial in $q^s+q^{-s}$.)
Note that these functions are traces of \emph{full} induced
representations: at a reducible point $\Theta_{I_{\overline{A}}(\chi)}$ is the sum
of the characters of the constituents, so the proposition does not by itself
describe the image of $\sF$ on irreducible tempered representations.
%\verify{whether the image of $C_c^\infty(\overline G)$ under the sum-of-constituents map is
%literally all of the stated space}

An analogous Schwartz-space image statement in the earlier manuscript is also recorded here for verification.

\begin{remark}[Schwartz principal-series image; verification pending]
For $f\in {\calC} (\overline{G })$, the proposed image
${\calC} (\overline{G })_{ps}^{\wedge }$
of the Fourier transform $f\longmapsto \phi_f$
with respect to the unitary principal series,
is given by 
\[
{\calC} (\overline{G })_{ps}^{\wedge }
= C^\infty(\Pi (\overline{C})_u)^{WH}.
\]
This identification requires the relevant Schwartz Paley--Wiener theorem
and topological control of all spectral components; it is not proved here.
\end{remark}


\subsection{The Fourier transform}
\label{subsec:FT}

Let
\[
\sJ={\rm span}\left\{f-f^y\ |\ f\in C_c^\infty(\overline{G })\ ,y\in 
\overline{G}\right\},
\]
where $f^y(x)=f(y^{-1}xy)$, and recall $PW_t(\overline{G })$ denotes the tempered Paley-Wiener space.

\begin{proposition}
\label{prop:kernel}
Let $\overline{G}$ be as in \S \ref{section:assume}. 
The kernel of the trace map
\[
\sF :C_c^\infty(\overline{G })\rightarrow PW_t(\overline{G }),
\]
defined by $\sF f(\pi)=\Theta_\pi(f)$, is $\sJ$.
\end{proposition}

\noindent\textbf{Proof and dependence on the literature.}
Every trace vanishes on $f-f^y$ because $\pi(f^y)$ is conjugate
to $\pi(f)$, so $\sJ\subset\ker\sF$.  For the converse, invoke
the tempered-character density theorem for covering groups
\cite[Theorem~5.8.10]{Li12}, which identifies the simultaneous kernel
of all genuine tempered character traces with the space generated
by conjugation differences.  Li's proof depends on the local invariant
trace formula and the representation-theoretic inputs described in
\cite[Sections~5.4--5.8]{Li12}; in particular his Theorem~5.4.1 is
admitted as a hypothesis for non-Archimedean fields.  The present
argument is therefore a deduction from that source under its stated
framework, not an independent derivation of the density theorem.
\hfill$\square$
\medskip

\begin{example}[A conjugation difference]\label{ex:difference}
If $f$ is any genuine compactly supported test function and $y\in\bG$,
then $h=f-f^y$ satisfies $\Theta_\pi(h)=0$ for every admissible
representation $\pi$.  The content of density is the converse for
all genuine tempered traces, not this elementary implication.
\end{example}



Let $V'$ denote the \emph{algebraic} dual of the complex vector space $V$.
If $V$ is in addition a $\overline{G }$-module, let 
$(V')^{\overline{G }}$ denote the
subspace of $\overline{G }$-invariant linear functionals.

\begin{lemma}
Let $\overline{G}$ be as in \S \ref{section:assume}. The canonical map
\[
(C_c^\infty(\overline{G })/\sJ)' 
\rightarrow (C_c^\infty(\overline{G })')^{\overline{G }}
\]
is an isomorphism.
\end{lemma}

\noindent
{\bf Proof.}
For $f\in C_c^\infty(\overline{G })$, let $f\ {\rm mod}\, \sJ$ 
denote its class in
$C_c^\infty(\overline{G })/\sJ$. First, note that the canonical map
\[
\begin{array}{c}
(C_c^\infty(\overline{G })/\sJ)'
\rightarrow (C_c^\infty(\overline{G })')^{\overline{G }}\\
D \longmapsto D^*\\
(f\ {\rm mod}\, \sJ\longmapsto D(f\ {\rm mod}\, \sJ)) \longmapsto
(f\longmapsto D(f\ {\rm mod}\, \sJ)).
\end{array}
\]
is injective by definition.

To see that this is surjective, let 
$D\in (C_c^\infty(\overline{G })')^{\overline{G }}$.
We must show that there is a $D_0\in (C_c^\infty(\overline{G })/\sJ)'$
such that $D=D_0^*$. Let
\[
D_0(f\ {\rm mod}\, \sJ)=D(f),\ \ \ \ \ \ f\in C_c^\infty(\overline{G }).
\]
We want to show that $D_0$ is well-defined,
i.e., that if $f,f'\in C_c^\infty(\overline{G })$ and
$f\ {\rm mod}\, \sJ=f'\ {\rm mod}\, \sJ$ then $D(f)=D(f')$.
By definition of $\sJ$, $f\ {\rm mod}\, \sJ=f'\ {\rm mod}\, \sJ$ implies
$f'=f+\sum_{i\in I}c_i(f_i-f_i^{y_i})$, for some finite set $I$ and
some $c_i\in {\C}$, $f_i\in C_c^\infty(\overline{G })$,
$y_i\in \overline{G }$. Since $D$ is invariant, from linearity it follows
that $D(f)=D(f')$, as desired. Therefore, the canonical map is 
surjective.
$\qed$

From these two results, we conclude the following
important fact.

\begin{theorem}
\label{thm:trace}
Let $\overline{G}$ be as in \S \ref{section:assume} above. The trace map $\sF$
induces an isomorphism $C_c^\infty(\overline{G})/\sJ\cong\sF(C_c^\infty(\overline{G}))\subset PW_t(\overline{G})$,
and hence, by the lemma, a canonical isomorphism
\[
\bigl(\sF(C_c^\infty(\overline{G}))\bigr)' 
\ \xrightarrow{\ \sim\ }\ (C_c^\infty(\overline{G })')^{\overline{G }}.
\]
In other words, every invariant distribution $D$ factors uniquely through the image of the tempered trace map. The duals in this theorem are algebraic duals; their use does not assert continuity of arbitrary factor functionals.  This is an algebraic statement about its values on compactly supported test functions.  It does not say that $D$ is a finite linear combination of tempered characters, or that its spectral functional is automatically continuous in a Schwartz Paley--Wiener topology.
\end{theorem}

This result allows us to define, for each invariant distribution $D$ on
$\overline{G}$, the {\it Fourier transform} $\widehat D$ on the algebraic image $\sF(C_c^\infty(\overline G))$
by
\begin{equation}
D(f)=\widehat D(\sF f),\ \ \ \ \ \ f\in C_c^\infty(\overline{G }).
\label{eqn:FT}
\end{equation}

By the results of \S \ref{sec:unitary}
above and of \cite{FK86},
the tempered dual has both a continuous
part and a discrete part. The continuous part of
$PW_t(\overline{G })$ decomposes
into a vector space sum of smooth functions on compact real tori.
Continuity on these parameter spaces is a separate topological assertion; it does not follow from the algebraic factorization alone.

\subsection{The Fourier transform as a spectral functional}
\label{subsec:main}

A function on the tempered dual extends linearly to the Grothendieck
group $R_t(\overline{G})$.  A general spectral distribution, however,
need not be determined by pointwise evaluations of such a function. 
The parameters of the Fourier transform are the \emph{irreducible} tempered
representations: if $I_{\overline{A}}(\chi)$ is reducible then $\Theta_{I_{\overline{A}}(\chi)}$ is only the sum of
the characters of its constituents, and a functional on $PW_t(\overline{G})$ may separate them.

The distributions that occur are of the following kind.

\begin{definition}\label{def:finite-order}
Let $X$ be a compact smooth torus.  A \emph{distribution of finite order}
on $X$ is a continuous linear functional on $C^\infty(X)$ for which,
for some $N$ and $C$, one has
\[
 |T(h)|\le C\sum_{|\alpha|\le N}\|\partial^\alpha h\|_\infty.
\]
For a spectral parameter space assembled from tori and discrete
representations, ``locally finite order'' means this property on each
of its smooth torus charts, subject to the compatibility at reducibility
points.  No uniform global bound over all Bernstein components is claimed.
\end{definition}

\begin{theorem}[Conditional spectral representation]\label{thm:main}
Assume the setting of Section~\ref{section:assume}.  Suppose the genuine
Schwartz trace transform extends to a continuous open surjection
\[
\sF:\calC(\bG)\longrightarrow PW_t^{\mathrm{Sch}}(\bG)
\]
onto a specified locally convex tempered Paley--Wiener space, with
$\ker\sF$ equal to the closure in $\calC(\bG)$ of the span of
conjugation differences.  Suppose also that the spectral space has local smooth-torus charts and
that each chart $X$ admits a continuous linear map
$E_X:C^\infty(X)\to PW_t^{\mathrm{Sch}}(\bG)$ realizing smooth
test functions on that chart, compatibly with its reducibility relations.
Then any invariant tempered distribution $D\in\calC(\bG)'$ induces
a unique continuous linear functional $\widehat D$ on
$PW_t^{\mathrm{Sch}}(\bG)$ satisfying
\begin{equation}\label{eq:conditional-pairing}
 D(f)=\langle\widehat D,\sF f\rangle,
 \qquad f\in\calC(\bG).
\end{equation}
On each smooth torus chart the induced spectral functional is a
distribution of finite order.  In general \eqref{eq:conditional-pairing}
is a \emph{distributional pairing}; no representation with an ordinary
scalar density $m(\pi)\,d\mu(\pi)$ is asserted.
\end{theorem}
\noindent\textbf{Proof.}
Because $D$ is invariant, it vanishes on every conjugation difference,
and by continuity it vanishes on their closure.  The kernel hypothesis
therefore gives $D|_{\ker\sF}=0$.  Define
$\widehat D(\sF f)=D(f)$.  This is well defined.  As $\sF$ is a
continuous open surjection, its quotient topology is the topology of
$PW_t^{\mathrm{Sch}}(\bG)$, so $\widehat D$ is continuous.  Surjectivity
gives uniqueness.  For a compact torus chart $X$, the composite functional
$\widehat D\circ E_X$ is continuous on $C^\infty(X)$ and is therefore
a distribution.  Every distribution on a compactsmooth manifold has finite order; Lemma~\ref{lemma:schwartz} proves
an explicit stronger representation for a torus.\hfill$\square$

\begin{remark}[Unverified analytic input]\label{rem:tpw-input}
The assumptions of Theorem~\ref{thm:main} are not proved in this paper.
Li \cite[Section~5.6]{Li12} establishes a tempered invariant
Paley--Wiener framework for covering groups.  Its precise topological
identification with the space and section conventions used in the
2001 argument requires comparison at reducibility parameters, including
the $R$-group hypothesis noted in Section~\ref{sec:lit}.
The conclusion is consequently recorded here as conditional rather
than as a new unconditional theorem.
\end{remark}


\medskip
The classical fact on which the local form rests is the following, which replaces the lemma
of L. Schwartz \cite{Sc50} (Ch.~III, Th.~XXI) quoted in the earlier version. (A distribution of finite order in the sense
above is locally $(1-\Delta)^Nh$, $h$ continuous.) (A single partial derivative of a
continuous function does not suffice: on a torus, $\langle\partial^Ih,1\rangle=0$ for $|I|\ge1$, while the Dirac mass
satisfies $\langle\delta_{x_0},1\rangle=1$.)

\begin{lemma} (Schwartz)
\label{lemma:schwartz}
Let $X=\R^d/(2\pi\Z)^d$ be a compact torus and $\Delta=\sum_j\partial^2/\partial x_j^2$. If $T$ is a distribution
on $X$ then there are an integer $N\ge0$ and a continuous function $h$ on $X$ such that
$T=(1-\Delta)^Nh$ (as distributions). If a finite group $W$ acts on $X$ by isometries and $T$ is $W$-invariant, $h$ may be chosen $W$-invariant.
\end{lemma}
\noindent
{\bf Proof.} Continuity of $T$ on $C^\infty(X)$ gives constants $c,\kappa$ with $|\widehat T(k)|\le c(1+|k|)^\kappa$ for
the Fourier coefficients $\widehat T(k)$, $k\in\Z^d$. Choose $N$ with $2N-\kappa>d$. Then
$h=\sum_k\widehat T(k)(1+|k|^2)^{-N}e^{ik\cdot x}$ converges absolutely to a continuous function and
$(1-\Delta)^Nh=T$ coefficientwise. Averaging $h$ over $W$ gives the invariant statement because the Laplacian commutes with the $W$-action. $\qed$



\subsection{Some corollaries}

By Theorem~\ref{thm:trace}, every invariant distribution on the compactly supported genuine test-function space factors algebraically through tempered traces.  Continuity of the spectral functional is an additional question.

\begin{remark}
Assume $\overline{G }=G$ ($n=1$).
In the earlier version it was claimed that if $D$ is tempered then it is admissible on the regular set,
that is, if we identify $C_c^\infty(G_{r})$ with the subspace of $C_c^\infty(G)$ of functions supported on
$G_r$, then $D|_{C_c^\infty(G_{r})}$ is admissible in the sense of \cite{HC78}. The proof in \cite{Jo01}
used the old form of \S\ref{subsec:main}; it should be re-derived from Theorem \ref{thm:main}.
%\verify{admissibility on the regular set}
\end{remark}

\begin{example}
 Clearly $f\longmapsto f(1)$ is an invariant tempered distribution.
Theorem \ref{thm:main} implies that its Fourier transform is a distribution of finite order on
$PW_t(\overline{G})$; in fact it is the Plancherel measure (order $0$):
\[
f(1)=\int_{\Pi(\overline{G })_t }
\Theta _\pi (f)\,d\mu_{\rm Pl}(\pi ),\;\;\;\;\;f\in 
C_c^\infty(\overline{G }). 
\]
This is a weak case of Harish-Chandra's Plancherel theorem.
\end{example}


\section{Summary of results and remaining problems}\label{sec:summary}
We distinguish results established here from those imported from the
literature or recorded conditionally.  The elementary factorization of
invariant linear functionals through a quotient is proved in
Section~\ref{subsec:FT}.  Proposition~\ref{prop:kernel}, the
substantive identification of that quotient with tempered traces,
is a consequence of a cited density theorem and its stated inputs
\cite[Theorem~5.8.10]{Li12}.  Theorem~\ref{thm:main} records a
conditional continuous Paley--Wiener formulation.  Lemma~\ref{lemma:schwartz}
is proved directly using Fourier series.

\begin{example}[An invariant distribution]\label{ex:identity}
Evaluation at the identity, $D(f)=f(1)$, is invariant under
conjugation.  It is tempered in the relevant Schwartz topology.
The genuine Plancherel formula identifies it, with the appropriate
normalization of measures, with a spectral functional on tempered
characters \cite[Section~2]{Li12}.  This is an example of a
spectral measure, unlike the derivative functional in
Example~\ref{ex:derivative}.
\end{example}

The following issues remain open \emph{in the present revision}:
(1) reconcile the opposite kernel--image identifications of the
endpoint intertwining matrix in the 1994 and 2026d manuscripts, and
verify the precise reducibility statements of Section~\ref{sec:unitary};
(2) check the proposed principal-series images against a full
Paley--Wiener description with finite Bernstein support and
reducibility relations;
(3) reconcile the topologies and $R$-group assumptions used for the
Schwartz extension with \cite{Li12}; and
(4) review all bibliographic page and theorem references against the
original sources.  These are tasks for an independent mathematical
referee, not claims of new open problems in representation theory.

\section*{Glossary}\label{sec:glossary}
\begin{description}[leftmargin=3.4cm,style=nextline]
\item[Genuine representation] A representation of $\bG$ on which
$(1,\zeta)$ acts by scalar $\zeta$ for every $\zeta\in\mu_n$.
\item[Genuine test function] A function satisfying $f((1,\zeta)x)=\zeta^{-1}f(x)$.
\item[Invariant distribution] A continuous linear functional $D$
with $D(f^y)=D(f)$ for all $y\in\bG$.
\item[Orbital integral] Integration of a test function over the
conjugacy orbit of a regular element, with a discriminant normalization.
\item[Tempered representation] An irreducible admissible representation
weakly contained in the regular representation, or equivalently described
by the usual tempered matrix-coefficient conditions in this setting.
\item[Paley--Wiener space] A space of trace transforms subject to
spectral regularity, support, and compatibility conditions.
\item[Schwartz space] Smooth functions with rapid decay relative to
Harish-Chandra's majorant $\Xi$ and the height $\sigma$.
\item[Cocenter] The quotient of the Hecke algebra by its subspace of
conjugation differences (equivalently by the commutator subspace).
\item[$R$-group] A finite group governing the constituents and
intertwining operators of certain tempered parabolic inductions.
\item[Distributional pairing] Evaluation of a continuous spectral
functional on a spectral test function, which need not be integration
against an ordinary function.
\end{description}

\section*{Acknowledgments}
The original 1994 notes and the article published in 2001 are by David
Joyner.  ChatGPT (OpenAI) contributed criticism of the 2000 text and
prepared this restructuring, mathematical status audit, and expository
revision in November 2026.  Claude (Anthropic) prepared the earlier
corrected 2026 draft, including its prior review notes.  The present
manuscript is an AI-assisted editorial and mathematical review;
\emph{AI review is not an independent certification of the results}.
The 2001 article's dedication to John Benedetto is retained in its
historical source rather than transferred to this differently framed
manuscript.


\begin{thebibliography}{99}

\bibitem{Ar80} H. Arit\"urk, 
 On the composition series of principal series
representations of a three-fold covering group of $SL(2,K)$,  
{\it Nagoya Math J.} {\bf 77}(1980)177-196 %\verify{initials}

\bibitem{Ar81} J. Arthur, 
 The trace formula in invariant form, {\it  Ann. Math.}
{\bf 114}(1981)1-74

\bibitem{Ar94} J. Arthur, 
 The trace Paley-Wiener theorem for Schwartz functions, 
{\it Contemp. Math.} {\bf 177}(1994)171-180

\bibitem{Be75} J. Benedetto, {\it Spectral Synthesis}, 
Academic Press, Inc., 1975

\bibitem{Be84} J. Bernstein, 
 Le center de Bernstein , (notes by P. Deligne), 
in {\it Repr\'esentations des 
groupes r\'eductifs sur un corp local},
Hermann, Paris, 1984

\bibitem{BZ76} J. Bernstein and A. Zelevinski,
 Representations of 
the group $GL(n,F)$, where $F$ is a non-archimedian local field, 
{\it Russian Math Surveys} {\bf 31}(1976)1-68

\bibitem{BZ77} J. Bernstein and A. Zelevinski,
 Induced representations of reductive $p$-adic 
groups, I,  {\it Ann. scient. \'Ec. norm. Sup.} {\bf 10}(1977)441-472

\bibitem{Da99} C. Davis, keynote talk presented at the
Conference on Harmonic Analysis in honor
of John Benedetto, University of Maryland,
October, 1999. 

\bibitem{FK86} Y. Flicker and D. Kazhdan, 
 Metaplectic correspondence,  
{\it Publ. Math. IHES} {\bf 64}(1986)53-110

\bibitem{Ge76} S. Gelbart, 
{\it Weil's Representation and the 
Spectrum of the Metaplectic Group},
Springer Lecture Notes in Math {\bf 530}, 1976

\bibitem{GS75} S. Gelbart and P. Sally, 
 Intertwining operators and 
automorphic froms on the metaplectic group,  
{\it Proc. Nat. Acad. Sci.}
{\bf 72}(1975)1406-1410

\bibitem{HC70} Harish-Chandra, {\it Harmonic analysis
on reductive $p$-adic groups},
(notes by G. van Dijk), 
Springer Lecture Notes in Math {\bf 162}, 1970

\bibitem{HC78} Harish-Chandra,
 Admissible invariant distributions on reductive $p$-adic groups, 
{\it Queen's papers in pure and applied math.} {\bf 48}(1978)
281-347

\bibitem{HC84} Harish-Chandra,
 The Plancherel formula for reductive $p$-adic groups ,
in {\it Collected Works}, vol IV, ed. V.S. Varadarajan,
Springer-Verlag, New York, 1984 (see also the paper
 Corrections ..., in the same volume)

\bibitem{Jo88}  D. Joyner,  On the metaplectic
analog of Kazhdan's `endoscopic' lifting,  
{\it Isr. J. Math.} {\bf 61}
(1988)113-154

\bibitem{Jo94}  D. Joyner,  A hitch-hiker's guide to genuine invariant 
distributions on metaplectic covers of SL(2) over a $p$-adic field, 
unpublished manuscript dated January 4, 1994 

\bibitem{Jo98}  D. Joyner,  A correspondence for the generalized Hecke 
algebra of
the metaplectic cover $\overline{SL(2,F)}$, $F$ $p$-adic, 
{\it New York J. of Math.} {\bf 4} (1998) 223-235. 

\bibitem{Jo99a}  D. Joyner,  On the unitary dual of 
$\overline{SL(2,F)}$, $F$ $p$-adic,  
preprint, 1999 

\bibitem{Jo99b} D. Joyner,  A correspondence for the supercuspidal
representations of $\overline{SL_2(F)}$, $F$ $p$-adic,  
{\it Archiv der Math.} 73(1999) 332--340.

\bibitem{Jo01}  D. Joyner,  Invariant distributions on the
$n$-fold metaplectic covers of $GL(r,F)$,
$F$ $p$-adic, \emph{J. Fourier Anal. Appl.} \textbf{7} (2001), 343--358.
(Original draft dated September 3, 2000)

\bibitem{JM00} D. Joyner and Courtney Moen, 
 The trace formula for $\overline{SL(2,\aaa)}$
 and applications , preprint, 2000.

\bibitem{Ka86} D. Kazhdan,  Cuspidal geometry of $p$-adic groups, 
{\it J. d'Anal. Math.} {\bf 47}(1986)1-36

\bibitem{KP84} D. Kazhdan and  S. Patterson, 
 Metaplectic forms,  {\it Publ. Math. IHES} {\bf 59}(1984)35-142

\bibitem{Li12} W.-W. Li, La formule des traces pour les rev\^etements de groupes r\'eductifs connexes.
II. Analyse harmonique locale, {\it Ann. Sci. \'Ec. Norm. Sup\'er.} {\bf 45}(2012)787-859.

\bibitem{Me98} P. Mezo,  A global comparison for general linear groups
and their metaplectic coverings,  Univ. of Toronto PhD thesis
(advisor: James Arthur), 1998

\bibitem{MoND} C. Moen,  Intertwining operators for covering groups 
of $SL(2)$ over a local field,  undated preprint (date unverified)

\bibitem{Mo88} C. Moen,  Irreducibility of unitary principal series 
representations for covering groups of $SL(2)$, 
{\it  Pac. J. Math.} {\bf 135}(1988)89-110

\bibitem{Cl26} Claude, notes prepared for D. Joyner, {\it Invariant distributions on metaplectic covers: a review
of Joyner's 2000 paper, a corrected theorem, and errata}, October 2026.

\bibitem{Sc50} L. Schwartz, 
{\it Theorie des distributions, vol. I}, Hermann, Paris, 1950

\bibitem{Si79} A. Silberger, {\it Introduction to
harmonic analysis on reductive $p$-adic groups}, 
Princeton Univ. Press, Princeton, 1979

\bibitem{Ta75} M. Taibleson, {\it Fourier analysis on local
fields}, Princeton Univ. Press, Princeton, 1975

\bibitem{Vi81} M.-F. Vign\'eras,  Caract\'erisation des integrales 
orbitales sur un groupe r\'eductif $p$-adique,  
{\it J. Fac. Sci. Univ. Tokyo }
{\bf 28}(1981)945-961

\bibitem{We75} A. Weil, {\it Basic Number Theory}, Springer-Verlag, 
1975

\end{thebibliography}



\end{document}